%%%%%%%%%%%%%%%%%%%%%%%%%%%%%%%%%%%%%%%%%%%%
%Set Document Type
%%%%%%%%%%%%%%%%%%%%%%%%%%%%%%%%%%%%%%%%%%%%
\documentclass[11pt,a4paper]{report}  

%%%%%%%%%%%%%%%%%%%%%%%%%%%%%%%%%%%%%%%%%%%%
%Packages
%%%%%%%%%%%%%%%%%%%%%%%%%%%%%%%%%%%%%%%%%%%%
\usepackage{MyStyle,graphicx}
\usepackage{xspace}
\usepackage[svgnames]{xcolor}
\usepackage{longtable}
\usepackage{amsmath}     %!  setspace is used to control linepacing
\usepackage[square]{natbib} %! needed for Harvard style of references.
\usepackage{enumerate}  %! used for the library form, but you might find 
\usepackage[hidelinks]{hyperref}
\usepackage[top=2.54cm, left=4cm, right=2.54cm]{geometry}
\usepackage{multicol}
\parskip
%%%%%%%%%%%%%%%%%%%%%%%%%%%%%%%%%%%%%%%%%%%%
%Set Tables
%%%%%%%%%%%%%%%%%%%%%%%%%%%%%%%%%%%%%%%%%%%%
\usepackage{lscape}
\usepackage{tablefootnote}
\usepackage{array}
\usepackage{bigfoot}
%%%%%%%%%%%%%%%%%%%%%%%%%%%%%%%%%%%%%%%%%%%%
%Set Font
%%%%%%%%%%%%%%%%%%%%%%%%%%%%%%%%%%%%%%%%%%%%
\usepackage[doublespacing]{setspace}
\usepackage[T1]{fontenc}
\usepackage{newpxtext,newpxmath}
%%%%%%%%%%%%%%%%%%%%%%%%%%%%%%%%%%%%%%%%%%%%
%Set Figure and Table Caption colour
%%%%%%%%%%%%%%%%%%%%%%%%%%%%%%%%%%%%%%%%%%%%
\usepackage{caption}
\captionsetup[table]{labelsep=period, labelfont={bf}, font={small, stretch=1.2}}     %% change 1.2 as you like
\captionsetup[figure]{labelsep=period, labelfont={bf}, font={small, stretch=1.2}}
%%%%%%%%%%%%%%%%%%%%%%%%%%%%%%%%%%%%%%%%%%%%
%Set Code Block Formatting
%%%%%%%%%%%%%%%%%%%%%%%%%%%%%%%%%%%%%%%%%%%%
\usepackage{listings}
\usepackage{color}
%%%%%%%%%%%%%%%%%%%%%%%%%%%%%%%%%%%%%%%%%%%%
%Set Table Block Formatting
%%%%%%%%%%%%%%%%%%%%%%%%%%%%%%%%%%%%%%%%%%%%
\usepackage{caption}
\captionsetup{width=0.98\textwidth}
%%%%%%%%%%%%%%%%%%%%%%%%%%%%%%%%%%%%%%%%%%%%
%Settings
%%%%%%%%%%%%%%%%%%%%%%%%%%%%%%%%%%%%%%%%%%%
%\thesisdraft                       %% Uncomment this if you want a draft
\leftchapter
\doublespacing
\usepackage{blindtext}
\usepackage{lastpage}
\usepackage{fancyhdr}
\pagestyle{fancy}
\fancyhf{} % clear existing header/footer entries
% Place Page X of Y on the right-hand
% side of the footer
\fancyfoot[R]{\small Page \thepage \hspace{1pt} of \pageref{LastPage}}
\renewcommand{\headrulewidth}{0pt}
%Prevent footnotes from breaking over multiple pages. 
\interfootnotelinepenalty=10000
%%%%%%%%%%%%%%%%%%%%%%%%%%%%%%%%%%%%%%%%                             
%Start Document
%%%%%%%%%%%%%%%%%%%%%%%%%%%%%%%%%%%%%%%%
\begin{document}

%%%%%%%%%%%%%%%%%%%%%%%%%%%%%%%%%%%%%%%%
%Title Pages
%%%%%%%%%%%%%%%%%%%%%%%%%%%%%%%%%%%%%%%%
\begin{titlepage}
    \begin{center}
        \singlespacing
        \Huge
        New Tools for the Identification of \\\vspace*{0.5cm} Fusarium Wilt in Banana \\
        \vspace*{1.75cm}
        \normalsize
        by \\ 
        \vspace*{0.2cm}
        \LARGE
        Jamie Pike BSc (Hons) \\ 
        \large
        \vfill
        \normalsize
        Submitted in the partial fulfilment of the requirements for the degree of \\
        \Large
        \vspace*{0.3cm}
        Doctor of Philosophy in Life Sciences \\
        \vspace*{1cm}
        \includegraphics[width=0.4\textwidth]{PrelimSections/crest_black.eps} \\
        \vspace*{1cm}
        \large
        University of Warwick, School of Life Sciences\\
        \vspace*{0.5cm}
        \normalsize
        January 2024 \\
         \vspace*{1cm}   
    \end{center}
\end{titlepage}

%%%%%%%%%%%%%%%%%%%%%%%%%%%%%%%%%%%%%%%%
%Fromatting
%%%%%%%%%%%%%%%%%%%%%%%%%%%%%%%%%%%%%%%%
\pagenumbering{roman} %! Begins roman numerals start from page i.
\clearpage
\tableofcontents                     
\listoftables                     
\listoffigures                    
%%%%%%%%%%%%%%%%%%%%%%%%%%%%%%%%%%%%%%%%
%Prelimary Sections
%%%%%%%%%%%%%%%%%%%%%%%%%%%%%%%%%%%%%%%%
\begin{thesisacknowledgments}        
\input PrelimSections/0.1_Acknowledgments.tex       
\end{thesisacknowledgments}

\begin{thesisdeclaration}       
\input PrelimSections/0.2_Declarations.tex 
\end{thesisdeclaration}

\begin{thesisabstract}              
  \begin{singlespace}    
    \input PrelimSections/0.3_Abstract.tex      
 \end{singlespace}
\end{thesisabstract}

\begin{thesisabbreviations}      
    \input PrelimSections/0.4_Abbreviations.tex     
\end{thesisabbreviations}

%%%%%%%%%%%%%%%%%%%%%%%%%%%%%%%%%%%%%%%%%%%%
%Main Chapters
%%%%%%%%%%%%%%%%%%%%%%%%%%%%%%%%%%%%%%%%
\pagenumbering{arabic} %! Begins arabic numerals start from page 1.
\renewcommand{\arraystretch}{1.5} %Extend tables so they're easier to read
%Set all instances of text below (e.g. \focub) to italicised full version.
\newcommand{\Focub}{\textit{Fusarium oxysporum} f. sp. \textit{cubense} }
\newcommand{\Fo}{\textit{Fusarium oxysporum }}
\newcommand{\Foxy}{\textit{Fusarium oxysporum }}
\newcommand{\Foc}{\textit{Foc }}
\newcommand{\et}{\textit{et al., }}
\newcommand{\Tef}\textit{Tef-1\(\alpha\) } 

%Set chapter names and .tex files with chapter content. 
\chapter{Introduction}
 \input Chapters/1_Introduction.tex
\chapter{Potential Novel \textit{Fusarium} Pathogen of Banana Identified }
 \input Chapters/2_PotentialNovelFusarium.tex
\chapter{Improving Genomic Tools for Identifying Virulence Factors in \textit{Fusarium oxysporum formae specialis}}
 \input Chapters/3_ImprovingGenomicTools.tex
\chapter{Phenomics: Image-based Analysis of Banana Stresses}
 \input Chapters/4_Phenomics.tex
 \chapter{Metabolomics: The Banana-pathogen Metabolome}
 \input Chapters/5_Metabolomics.tex


%%%%%%%%%%%%%%%%%%%%%%%%%%%%%%%%%%%%%%%%%%%%
%Appendices
%%%%%%%%%%%%%%%%%%%%%%%%%%%%%%%%%%%%%%%%%%%%
\appendix                            %% this will do the appendices
\chapter{Scripts used for \textit{Fusarium} Genome Assembly and Analysis}
For all of the scripts written as part of the genome assembly and analysis of the TNAU \textit{Fusarium} isolates, SY-2, S6, S16, and S32. 
%%%%%%%%%%%%%%%%%%%%%%%%%%%%%%%%%%%%%%%%%%%%%%%%
%CODE OUTPUT STYLE
%%%%%%%%%%%%%%%%%%%%%%%%%%%%%%%%%%%%%%%%%%%%%%%%

\lstset{basicstyle=\ttfamily,
  showstringspaces=false,
  commentstyle=\color{red},
  keywordstyle=\color{blue}
}

\lstdefinestyle{customc}{
  belowcaptionskip=1\baselineskip,
  breaklines=true,
  showstringspaces=false,
  basicstyle=\footnotesize\ttfamily,
  keywordstyle=\bfseries\color{green!40!black},
  commentstyle=\itshape\color{purple!40!black},
  identifierstyle=\color{blue},
  stringstyle=\color{orange},
}

%Stops the script running off the page
\lstset{escapechar=@,style=customc}

%%%%%%%%%%%%%%%%%%%%%%%%%%%%%%%%%%%%%%%%%%%%%%%%
%CODE SECTIONS
%%%%%%%%%%%%%%%%%%%%%%%%%%%%%%%%%%%%%%%%%%%%%%%%

\section{GCTrimmer.py}\label{apx:gcTrimmer.py}
Script used to separate contigs in the TNAU assemblies into separate FASTA files based on GC\% threshold. 

\begin{lstlisting}[language=Python, caption=Script to separate contigs based on GC\% threshold]
#Extracts sequences from a FASTA file and order in to two different FASTAs depending on GC content

###############################################
#Set up and import modules.
import re, sys, colorama
from datetime import date
from tqdm import tqdm
from colorama import Fore
from Bio import SeqIO, SeqUtils
from Bio.SeqUtils import GC
from Bio.SeqRecord import SeqRecord
###############################################
#Establish input file (FASTA) and gc Threshold.
infile = sys.argv[1]
gcContentThreshold = sys.argv[2]
###############################################
#Establish colour reset.
colorama.init(autoreset=True)
###############################################
# Establish Count for sequences above and below threshold.
SeqAboveCount = 0
SeqBelowCount = 0

###############################################
#Progress:
#Record date
today = date.today()
startTime = today.strftime("%B %d, %Y")
#Create progress bar
num = len([1 for line in open(infile) if line.startswith(">")])
print("="*80)
print(f'Running gcTrimmer.py.\n{startTime}.\n\nInput File:\t\t{infile}\nNumber of Contigs:\t{num}\nGC threshold:\t\t{gcContentThreshold}%\n')

#Start progress bar
with tqdm(total=num) as pbar:
###############################################
    #Open two .fasta files using the name of the infile.
    with open(f'{infile}_GCcontentBelow{gcContentThreshold}perc.fasta'.format(), "a") as LessThanFile, open(f'{infile}_GCcontentAbove{gcContentThreshold}perc.fasta'.format(), "a") as GreaterThanFile, open(f'gcContentReport.txt'.format(), "a") as report  :

    #Parse and loop through the FASTA input file. Search each scaffold/contig for sequences with length => than the desired GC content and deposit them in the FASTA
    #containing the less than the desired gc content. If the sequence gc content is greater than the gc content threshold, then it is added to the "greater than" FASTA file.
        for sequence in SeqIO.parse(infile, "fasta"):
            print("="*10, file=report, sep="\n")
            gccontent = (sequence.seq.count('G') + sequence.seq.count('C')) / len(sequence) * 100
            pbar.update(1)
            if gccontent < int(gcContentThreshold):
                print(f'{sequence.description} is below the threshold GC content. GC content is {Fore.RED}{gccontent}{Fore.BLACK}.\n', file=report, sep="\n")
                print(sequence.format("fasta"), file=LessThanFile, sep="\n")
                SeqBelowCount = SeqBelowCount + 1
            else:
                print(f'{sequence.description} is above the threshold GC content. GC content is {Fore.GREEN}{gccontent}{Fore.BLACK}.\n', file=report, sep="\n")
                print(sequence.format("fasta"), file=GreaterThanFile, sep="\n")
                SeqAboveCount = SeqAboveCount + 1


#Print the total number of sequences above and below the GC threshold.

print(f'\nThe number of sequences above the GC threshold: {Fore.GREEN}{SeqAboveCount}{Fore.BLACK}.\nThe number of sequences below the GC threshold: {Fore.RED}{SeqBelowCount}{Fore.BLACK}.')
print("="*80)    
\end{lstlisting}

\section{ContaminatFilter.sh}\label{apx:ContamFilter}
\begin{lstlisting}[language=Python, caption=bash script to filter contigs from assembly that are predicted to belong to a non-\textit{Fusarium} genus by Blobtools.]
#!/bin/bash

#Jamie Pike
#Command for filtering conatminant blobtools hits. This groups all of the standard BlobTools commands I used to remove a contaminat contigs analysis. 

python -c "print('=' * 75)"
echo "Blobtools Contaminant Filter"
echo "----------------------------"
echo $(date)
echo "Usage: ContaminantFilter.sh <FASTA file> <blobtools json file> <outfile prefix>"
python -c "print('=' * 75)"

infile=${1?Please provide the assembly input fasta.} #Input Assembly.
json=${2?Please provide a BlobTools json file.} #Input BAM.
prefix=${3?Please provide a prefix for you outputs.} 

########################
#Escape the script if there are any errors. 
set -e 

echo "Creating species table..."
#Use BlobTools view to generate a blobtools table.txt file for filtering.
blobtools view -i ${json} -r species -o ${prefix}

echo "Filetering for Fusarium and no-hit only contigs..."
#Use grep to extact only the Fusarium and no-hit lines from the hit table.
grep -E 'Fusarium|no-hit' ${prefix}*.table.txt > ${prefix}.FusariumHits.table.txt

echo "Generating list file..."
#Create a list of the Fusarium and no-hit only contigs. 
awk '{print $1}' ${prefix}.FusariumHits.table.txt >  ${prefix}.FusariumHits.list.txt

echo "Generating contaminant filtered fasta..."
#Create a list of the Fusarium and no-hit only contigs. 
blobtools seqfilter -i ${infile} -l ${prefix}.FusariumHits.list.txt -o ${prefix}.ContaminantFiltered

echo "Generating contaminant sequences fasta..."
#Create a list of the Fusarium and no-hit only contigs. 
blobtools seqfilter -v -i ${infile} -l ${prefix}.FusariumHits.list.txt -o ${prefix}.ContaminantSequences

echo "ContaminantFilter.sh finished."
echo $(date)
echo "+++++
It is advisable that you check the number of contigs in the ${prefix}.ContaminantFiltered.fasta matches the number of lines in the ${prefix}.FusariumHits.table.txt file."
python -c "print('=' * 75)" 

\end{lstlisting}



\chapter{Scripts used as part of Maei Pipeline and Subsequent Analysis}
For all of the scripts written for the \textit{mimp}-associated effector identification pipeline, as well as scripts used for the additional analysis of the \textit{Fusarium} database used in the identification of effector candidates. 
%%%%%%%%%%%%%%%%%%%%%%%%%%%%%%%%%%%%%%%%%%%%%%%%
%CODE OUTPUT STYLE
%%%%%%%%%%%%%%%%%%%%%%%%%%%%%%%%%%%%%%%%%%%%%%%%

\lstset{basicstyle=\ttfamily,
  showstringspaces=false,
  commentstyle=\color{red},
  keywordstyle=\color{blue}
}

\lstdefinestyle{customc}{
  belowcaptionskip=1\baselineskip,
  breaklines=true,
  showstringspaces=false,
  basicstyle=\footnotesize\ttfamily,
  keywordstyle=\bfseries\color{green!40!black},
  commentstyle=\itshape\color{purple!40!black},
  identifierstyle=\color{blue},
  stringstyle=\color{orange},
}

%Stops the script running off the page
\lstset{escapechar=@,style=customc}

%%%%%%%%%%%%%%%%%%%%%%%%%%%%%%%%%%%%%%%%%%%%%%%%
%CODE SECTIONS
%%%%%%%%%%%%%%%%%%%%%%%%%%%%%%%%%%%%%%%%%%%%%%%%

\section{MimpFinditer.py}\label{apx:mimpfinditer}

The python script below uses the \textit{mimp} TIRs to identify \textit{mimps} in \textit{Fusarium oxysporum} genomes, recording their locations in a .bed file. 

\begin{lstlisting}[language=Python, caption= Regex based \textit{mimp} searching script.]
#!/bin/python

#### MIMP FINDER - SEARCHES FOR ALL MIMP IN SEQUENCE ####

#Search through a FASTA file and generate a bed file containing the location of mimps.

#Set up and import correct modules.

from datetime import datetime
startTime = datetime.now()
import re, sys

from Bio import SeqIO
# from Bio.Seq import Seq
from Bio.SeqRecord import SeqRecord

############################################
#Establish input file (FASTA) and start time.

infile = sys.argv[1]

print(f'\n\nRunning Mimp_finder.py on {infile} at {startTime}.')

#Open a .bed file using the name of the infile. The loctation of the mimp finder will be desposited within.

with open(f'{infile}_mimp_regex_hits.bed'.format(), "a") as file_object:

#Parse and loop through the FASTA input file. Search through each scaffold/contig of the FASTA looking for sequences which match the regex/mimp TIR "CAGTGGG..GCAA[TA]AA" (re.IGNORECASE remove case sensitivity issues).
#mimp_length = Create a variable for mimp length - Mimps cannot exceed 400 nucleotides.
#mimp_count = Create a variable to count the number of mimps found on each contid/scaffold.

    for sequence in SeqIO.parse(infile, "fasta"):
        print("="*80 + "\n")
        print(f'Searching {sequence.description} for mimps.\n')
        hit = re.finditer(r"CAGTGGG..GCAA[TA]AA", str(sequence.seq), re.IGNORECASE)
        mimp_length = 400
        mimp_reverse_length = -400
        mimp_count = 0

#For every instance where the mimp TIR "CAGTGGG..GCAA[TA]AA" is found, print the location of the hit and search that contig for the other mimp TIR.
#h_start = record the start location of the mimp

        for h in hit:
            if h:
                h_start = h.start()
                print(f'--Mimp TIR found: "{h.group()}" at position {h.start()} to {h.end()}--')
                hit_rc = re.finditer(r"TT[TA]TTGC..CCCACTG", str(sequence.seq), re.IGNORECASE)

#Ensure that the distance between the start of the first mimp TIR and end of the second mimp TIR is not greater than mimp_length (400 nucleotides).
#If TRUE, add 1 to the mimp_count variable and print the contig/scaffold description along with the start location of the mimp and the end location of the mimp in the .bed file opened (ensure fields are \t seperated to maintain .bed file convention).
#If the distance between the start of the first mimp TIR and the end of the second mimp TIR is greater than 400 nucleotides, it is not a mimp.
#If the second mimp TIR comes before the first mimp TIR found (second mimp TIR end - first mimp TIR start < 0), it is not a mimp.

                for h_rc in hit_rc:
                    print('Looking for reverse complement sequence(s)...')
                    h_rc_end = h_rc.end()
                    if h_rc:
                        print(f'Mimp reverse complement TIR found: "{h_rc.group()}" at position {h_rc.start()} to {h_rc.end()}.  Is this complete mimp?')
                        length = h_rc_end - h_start
                        if length > mimp_reverse_length:
                            if length < mimp_length:
                                print(f'Yes! Full mimp found at position {h.start()} to {h_rc.end()}.')
                                print("The mimp length is:")
                                print(f'{length}bp\n')
                                mimp_count = mimp_count+1
                                print(sequence.description, h.start(), h_rc.end(), "+", file=file_object, sep="\t")
                            else:
                                print("No, TIRs are at a greater distance than 400bp. \nOnly a partial mimp.\n")
                        else:
                            print("No, TIRs are at a greater distance than 400bp")

#Print the number of mimps found on that contig/scaffold.
#When whole FASTA has been looped through, print the file searched and time finished.


                print(f'\n{sequence.description} contains {mimp_count} mimp(s)')


    print("="*80 + "\n")

    print(f'\nMimp_finder.py complete. \n\nFile searched: {infile}\n\nFinished at: {startTime}')
\end{lstlisting}


%\bibliographystyle{harvard}
%\bibliography{Bibliogrphy}            %% Start your bibliography here;

                                 %! with sample.bib as your bibliography file. You can
                               %% also use:
                %\begin{thebibliography}
                 %\bibitem{Bibliography/NewTools Bib.bib}
                %\end{thebibliography}
                               %% to generate your bibliography.

%\begin{thesisauthorvita}             %% Write your vita here; it can be
%                                     %% anything in LaTeX2e par-mode.
%\end{thesisauthorvita}               %%




\end{document}                       %% Done.

